\subsection{Diferencias entre clases de reacción registradas y comportamiento químico}
El conjunto reúne tanto rutas en las que el compuesto de interés se transforma en otro producto como rutas en las que dicho compuesto es el producto buscado. Separar estos roles es esencial: una hidrogenación registrada para obtener un alcano no equivale a la hidrogenación del propio alcano como reactante. La asimetría observada en la Tabla~\ref{tab:clases} no demuestra una propiedad intrínseca de la sustancia, porque el conjunto se construyó deliberadamente con diez ejemplos por rol y compuesto.

Los alquenos poseen un enlace $\pi$ en el doble enlace C=C, circunstancia que ayuda a comprender la relevancia de transformaciones de adición. Los alcanos, al carecer de dobles enlaces C=C, presentan rutas y requisitos distintos. No obstante, la clasificación textual de los registros no permite afirmar, por sí sola, el mecanismo de cada transformación. En consecuencia, categorías amplias como «otro», «reducción» y «alquilación» se describen aquí como etiquetas documentales y no como mecanismos químicos establecidos para cada registro.

\subsection{Estructura y propiedades físicas: el caso de los isómeros}
La diferencia de 11,2~$^\circ$C entre butano e isobutano del archivo ofrece una comparación especialmente informativa porque ambas moléculas tienen igual fórmula molecular y masa molar. La ramificación altera la geometría y la extensión del contacto intermolecular, lo que puede modificar la magnitud promedio de las interacciones dispersivas y el punto de ebullición. El isobutano es el isómero ramificado y presenta el punto de ebullición más bajo entre ambos, en concordancia con la explicación cualitativa esperada. Sin embargo, esta comparación de dos compuestos no constituye una prueba estadística independiente de causalidad.

En contraste, comparar indiscriminadamente densidades, solubilidades o presiones de vapor de todas las sustancias relacionadas puede resultar engañoso. La fuente mezcla medidas reportadas a distintas temperaturas, presiones y posibles fases físicas. Por ejemplo, una densidad informada para una sustancia gaseosa no puede interpretarse sin más como si hubiera sido medida en las mismas condiciones que la de un líquido. Por ello, las gráficas de dichas magnitudes se mantienen como exploratorias y no sirven aquí para sostener correlaciones generales.

\subsection{Condiciones térmicas y alcance de la comparación}
La mediana de 100~$^\circ$C caracteriza el conjunto de temperaturas centrales extraídas, no la demanda energética de un tipo de reacción. La reacción puede transcurrir en varias etapas térmicas, a presión elevada o bajo catalizadores distintos. Incluso una misma clase textual puede abarcar procesos experimentalmente heterogéneos. Cualquier comparación rigurosa de cinética, selectividad o rendimiento requeriría variables adicionales no consolidadas en esta base y una selección de experimentos bajo condiciones comparables.

\subsection{Calidad de datos y limitaciones}
Las verificaciones automáticas ofrecen evidencia de correspondencia documental, pero no sustituyen la validación de la fuente química. Destacan las siguientes limitaciones: (i) la muestra no fue seleccionada de forma aleatoria; (ii) algunos identificadores de reacción reaparecen asociados a varios compuestos; (iii) las propiedades de productos o reactantes pueden repetirse cuando una misma sustancia figura en diferentes registros; (iv) hay numerosos faltantes y datos cualitativos; (v) la normalización textual puede reunir bajo una categoría rutas distintas; y (vi) ciertas propiedades carecen de condiciones de medida homogéneas. Es especialmente importante no usar valores faltantes como ceros ni tratar repeticiones como moléculas independientes en una correlación.

La auditoría no detectó diferencias textuales críticas en los campos cotejados, y las 1.060 conversiones automáticamente comprobables concordaron. El resultado respalda la consistencia de la transferencia y estandarización de los valores verificables. Las otras 620 celdas no aportan evidencia de discrepancia: representan casos fuera del alcance de esas reglas, como datos ausentes, cualitativos o expresados bajo condiciones heterogéneas.
